% Agent 15: Pages 399–402 (PDF pp.28–31)
% Sections 4.5 Interaction of Outflowing Radiation with Gas
%          4.6 Validity of Hydrodynamical Approximation
%          4.7 Gas Flow Near the Horizon
% Equations (4.5.1) through (4.7.7)

\documentclass[aps,prd,twocolumn,superscriptaddress,nofootinbib]{revtex4-2}
\usepackage{amsmath,amssymb,amsfonts}
\usepackage{graphicx}
\usepackage{bm}
\usepackage{mathrsfs}
\usepackage{hyperref}
\usepackage{physics}

% Custom commands
\newcommand{\fvec}[1]{\mathbf{#1}}

\begin{document}

%% ============================================================
%% Page 399 — continuation from S4.4 and start of S4.5
%% ============================================================

$\nu \sim 7 \times 10^{14}\;\text{Hz}$
($\lambda \sim 4000\;\text{\AA}$); and it will die out exponentially for
$\nu > 7 \times 10^{14}\;\text{Hz}$; cf.\ \S2.3.
As \cite{Schwarzman1971} remarks, such a spectrum
and luminosity resemble those of ``DC white-dwarf stars'' (white dwarfs with
continuous, line-free spectra).  This has led Schwarzman to speculate that
``perhaps some of the objects heretofore regarded as type DC white dwarfs are
actually black holes.''

Instabilities and inhomogeneities in the flow (due, e.g., to reconnection of
field lines) might lead to marked fluctuations in the luminosity of the hole.  The
timescales for such fluctuations might be of the order of the gas travel time from
${\sim}10\,r_g$ to ${\sim}2r_g$; i.e.

\begin{equation}
\Delta t_{\text{fluctuations}} \sim (10^{-3}~\text{to}~10^{-4}~\text{sec})(M/M_\odot)\,.
\tag{4.4.8}
\end{equation}

Such rapid fluctuations in an object so faint should be very difficult to detect,
even with the 200-inch telescope.  For further details on the fluctuations, see
the end of \S4.7.

%% ============================================================
%% S4.5 — Interaction of Outflowing Radiation with the Gas
%% ============================================================

\subsection{Interaction of Outflowing Radiation with the Gas}
\label{sec:4.5}

As the luminosity $L$ pours out from the region $r \sim 2r_g$, it interacts with the
inflowing gas.  For the luminosities $L \sim 10^{32}$~ergs/sec and frequencies $\nu \sim 10^{14}$~Hz
derived in the last section, one can readily verify that the infalling gas is optically
thin to free-free absorption and to synchrotron reabsorption.  (See \S2.6 for
basic concepts and equations used in such a calculation.)

What of the electron concentration is sufficiently high, it can act as a ``blanket'' to
impede the outflow of the radiation.  To test for such blanketing one can examine
the optical depth of the gas to electron scattering.

\begin{equation}
\tau_{es}(r = 2r_g) = \int_{2r_g}^{\infty} \kappa_{es} \rho\, dr
\simeq (1 \times 10^{-9})(M/M_\odot)(\rho_\infty/10^{-24}~\text{g\,cm}^{-3})(T_\infty/10^4\,\text{K})^{-3/2}\,.
\tag{4.5.1}
\end{equation}

(A more careful calculation would replace the nonrelativistic absorption
coefficient $\kappa_{es} = 0.4$~cm$^2$/g by the coefficient appropriate to relativistic
electrons, since $T \sim 10^{12}$~K at $r \sim r_g$; such a calculation would also reveal
extreme optical thinness.)  (ii) Since the outgoing photons have energies
$h\nu \sim 10$~eV far less than the electron kinetic energies, Compton scattering can
modify the spectrum, producing high-energy photons [see eq.\ (2.4.9)].  However,
the extreme optical thinness guarantees that only about one photon in a million
scatters; so this effect can be ignored.  (iii) The outpouring photons exert a
pressure on the infalling gas when they scatter, and this effect can retard the
inflow.  Notice that the magnitude of the first two effects is governed by the

\bigskip
%% ============================================================
%% Page 400
%% ============================================================

density of the inflowing gas.  By contrast, the magnitude of the pressure effect
is governed by the luminosity of the outpouring radiation.  Let us calculate this
effect explicitly, at radii sufficiently large that the electrons are nonrelativistic.

Because the electron scattering cross-section $d\sigma/d\Omega = \tfrac{1}{2}r_0^2(1+\cos^2\theta)$ is
symmetric between forward and backward angles, $h\nu/c$, to the electron.  Hence, the time-averaged
photon force acting on an electron at radius $r$ is

\begin{equation}
\langle\text{Force}\rangle = \left\langle\frac{d(\text{momentum})}{dt}\right\rangle
= \int \frac{d(\text{number of photons})}{d(\text{area})\,dt\,d\nu}\, h\nu\,\sigma_e\, d\nu
\tag{4.5.2}
\end{equation}
%
\[
= \sigma_{eF} F = \sigma_e (L/4\pi r^2)\,,
\]
%
where $\sigma_e = (8\pi/3)r_0^2 = 0.657 \times 10^{-24}$~cm$^2$ is the total scattering cross section,
$F$ is the radiation flux (ergs~cm$^{-2}$~sec$^{-1}$) and $L$ is the total luminosity of the hole.
Notice that this time-averaged force is completely independent of the spectrum
of the radiation (so long as most of the photons have $h\nu \ll m_e c^2$ as seen in the
electron rest frame, so that nonrelativistic cross sections are applicable).  It is
also independent of the optical thickness---being equally valid for $\tau \gg 1$, in
which case

$F \ll J$ (see \S2.6 for notation).

The pull of gravity must work against this $\sigma_e^{} (L/4\pi r^2)$ force.  The photon force
acts almost entirely on the electrons ($\sigma_e \sim 1$~(mass)$^{-2}$, so each proton feels a
force $3 \times 10^6$ times weaker than that felt by an electron).  By contrast, the
acceleration of gravity acts equally on electrons and photons.  Hence, the electrons
move outward and slightly relative to the photons, creating an electric field that
transmits the photon force $\sigma_e L/4\pi r^2$ from electrons to protons.  The net force
that then acts on proton (assuming a completely ionized hydrogen gas) is

\begin{equation}
\langle\text{Force}\rangle_{\text{total}} = \frac{\sigma_e L}{4\pi r^2} - \frac{GMm_p}{r^2}\,.
\tag{4.5.3}
\end{equation}

Thus, there is a critical luminosity [``Eddington'' (1926) limit]

\begin{equation}
L_{\text{crit}} \equiv \frac{4\pi G M m_p}{\sigma_e}
= (1.3 \times 10^{38}~\text{ergs/sec})\,(M/M_\odot)
\tag{4.5.4}
\end{equation}

such that for $L \ll L_{\text{crit}}$ gravity dominates at all radii, but for $L \gg L_{\text{crit}}$ photon
pressure dominates at all radii.  \cite{Zeldovich1964} derived this limit for the case of
stellar equilibrium.  \cite{Zeldovich1964} derived this limit for the case of
theory for the cases of quasars and accretion.

For the case of current interest---a black hole of $M\sim 1$ to $100\,M_\odot$ swallowing
interstellar gas---the luminosity is far below the critical value, so the photon
pressure can be ignored.

%% ============================================================
%% Page 401
%% ============================================================

However, in other contexts the accretion rate $\dot{M}_0$ may be so great, and the
efficiency $\xi$ for converting the inflowing mass-energy into outgrowing luminosity
$L$ may be sufficiently high that

\begin{equation}
\frac{L}{L_{\text{crit}}} = \frac{\xi \dot{M}_0 c^2}{L_{\text{crit}}}
\lesssim \frac{\xi}{0.1} \left(\frac{\dot{M}_0}{10^{18}~\text{g/sec}}\right)
\tag{4.5.5}
\end{equation}

may temporarily exceed unity.  When this happens, the photon pressure will
impede further mass inflow, thereby reducing the luminosity to $L \lesssim L_{\text{crit}}$.  Such
accretion is said to be ``\emph{self-regulated}''.  Self-regulated accretion, with modifications due to lack of spherical symmetry, may be important for black holes in
close binary systems; see \S5.13.  However, it is unimportant for the case at
hand.

The outflowing radiation from a black hole in interstellar space ($L \sim 10^{32}$
ergs/sec, $\nu \sim 10^{14}$--$10^{15}$~Hz) can also interact with the surrounding interstellar
gas.  One can verify that the radiation is sufficiently intense to keep the gas at a
temperature $T_\infty \sim 10^4$~K in the neighborhood of $r = r_t$, and to keep it partially
ionized.

%% ============================================================
%% S4.6 — Validity of the Hydrodynamical Approximation
%% ============================================================

\subsection{Validity of the Hydrodynamical Approximation}
\label{sec:4.6}

The above model for accretion onto an interstellar black hole relies crucially on
the validity of the hydrodynamical approximation at and outside the sonic point,
and on equipartition assumptions inside the sonic point.  If the gas were to behave
like independent particles rather than like a fluid, then the accretion rate would
be greatly reduced (see \S4.1).  Thus, it is crucial to check the validity of the
hydrodynamical approximation.

In an ionized hydrogen plasma \emph{without} magnetic fields, the deflection of
protons away from straight-line paths is caused primarily by the cumulative
effects of many small-angle Coulomb scatterings.  The total distance a proton
must travel before the ``random-walk'' deflections add up to an angle of ${\sim}90^\circ$
is [see, e.g., \cite{Shkarofsky1966} or \cite{Pikelner1961}]

\begin{equation}
\lambda_p = \frac{9}{4\pi}\,\frac{(kT)^2}{n_p e^4 L_c}
\simeq (7 \times 10^{12}~\text{cm})\left(\frac{\rho_\infty}{10^{-24}~\text{g\,cm}^{-3}}\right)^{-1}
\left(\frac{T}{10^4\,\text{K}}\right)^{2}\,.
\tag{4.6.1}
\end{equation}

Here $L_c$ is a ``Coulomb-logarithm'' (analogue of Gaunt factor in theory of
Bremsstrahlung) given by

\begin{equation}
L_c \simeq 23 + \tfrac{3}{2}\ln (n_e/n_{\text{cm}^{-3}});
\tag{4.6.2}
\end{equation}

$n_e$ and $n_p$ are the number densities of electrons and (ionized) protons; and we
have evaluated $\lambda_p$ for the case of a fully ionized hydrogen plasma.  The plasma
will behave like a fluid on scales $l \gg \lambda_p$, but like independent particles on scales
$l \ll \lambda_p$.

The scale of interest for determining the accretion rate is the sonic radius

\begin{equation}
r_s = \frac{5 - 3\Gamma}{4}\,\frac{GM}{a_\infty^2}
\simeq (10^{13}~\text{cm})\left(\frac{M}{M_\odot}\right)
\left(\frac{T_\infty}{10^4\,\text{K}}\right)^{-1}
\tag{4.6.3}
\end{equation}

$(\ll \lambda_p)$.  The scale of interest for $\left(\frac{M}{M_\odot}\right)\left(\frac{T_\infty}{10^4\,\text{K}}\right)^{-1}$

at which point $T \sim T_\infty$, $\rho \sim \rho_\infty$, so

\begin{equation}
\lambda_p(r_s) \sim 0.7 \left(\frac{M}{M_\odot}\right)^{-1}
\left(\frac{T_\infty}{10^4\,\text{K}}\right)^{-1}\,.
\tag{4.6.4}
\end{equation}

%% ============================================================
%% Page 402
%% ============================================================

Thus, if one ignores the magnetic field, then the gas will not quite behave like a
fluid near the sonic point---and it may fail even more to be fluid-like at small
radii; cf.\ \S\,13.4 of ZN.

The interstellar magnetic field, of strength $B_\infty \sim 10^{-6}$~G, can ``save'' the
hydrodynamic approximation.  It deflects protons away from straight-line
motion in a distance of the order of the Larmour radius

\begin{equation}
\lambda_L = \frac{m_p c\,v}{eB}
\simeq (1 \times 10^{8}~\text{cm})
\left(\frac{B}{10^{-6}~\text{G}}\right)^{-1}
\left(\frac{T}{10^4\,\text{K}}\right)^{1/2}
\tag{4.6.5}
\end{equation}

and this is true even when the thermal pressure of the plasma exceeds the
magnetic pressure so that the field gets dragged about by the gas.  The Larmour
radius is small compared to all macroscopic scales of interest (e.g.\ small compared
to $r_S$ when one uses the values $B \sim B_\infty$ and $T \sim T_\infty$ appropriate to $r_S$).  Hence,
the hydrodynamical approximation is fairly well justified.

%% ============================================================
%% S4.7 — Gas Flow Near the Horizon
%% ============================================================

\subsection{Gas Flow Near the Horizon}
\label{sec:4.7}

Near the horizon, $r \sim r_g$, relativistic effects will modify the Newtonian picture
of accretion.  It is in fact difficult to solve for the relativistic flow using the
relativistic Bernoulli equation (2.5.33) and the law of mass conservation (2.5.23').
Such a solution is analogous to the Newtonian hydrodynamical solution studied
in \S4.2.  However, such a solution is partially irrelevant because it ignores the
role of the magnetic field, and such a solution is not needed because the intensity
of the gravitational pull, as viewed in stationary frames near the horizon, is so
great as to guarantee supersonic flow with qualitatively the same form as free-particle
fall.  Thus, from a knowledge of geodesic orbits one can infer the main
features of the flow.

Consider, first, spherical accretion onto a nonrotating hole (Schwarzschild
gravitational field).  A fluid element in the accreting gas, like a freely falling
particle, reaches the horizon in finite proper time.  The moment of crossing the
horizon is in no way peculiar, as seen by proper time.  The density does not
reach infinity there; and the temperature does not reach infinity.  In fact,
because the equation for radial geodesics in the Schwarzschild geometry

\begin{equation}
ds^2 = (1 - r_g/r)\,c^2\,dt^2 + (1 - r_g/r)^{-1}\,dr^2 + r^2\,d\Omega^2\,.
\tag{4.7.1}
\end{equation}

has the first integral

\begin{equation}
(dr/d\tau)^2 = 2GM/r \quad\text{(for fall from rest at } r \gg r_g\text{)},
\tag{4.7.2}
\end{equation}

and because the law of rest-mass conservation $\boldsymbol{\nabla}\cdot(\rho \fvec{u}) = 0$ has the first integral

\begin{equation}
4\pi r^2 \sqrt{-g}\, \rho\,(dr/d\tau) = \dot{M}_0\,,
\tag{4.7.3}
\end{equation}

the density of rest mass must vary as

\begin{equation}
\rho \simeq \left(6 \times 10^{-12}~\text{g\,cm}^{-3}\right)
\left(\frac{M}{M_\odot}\right)^{-1}
\left(\frac{T_\infty}{10^4\,\text{K}}\right)^{-3/2}
\,(r/r_g)^{-3/2}\,.
\tag{4.7.4}
\end{equation}

This must be as true near and inside the horizon, $r \lesssim r_g$, as it is in the Newtonian
realm $r \gg r_g$.  Similarly, the temperature will retain its Newtonian form (4.4.2):

\begin{equation}
T \simeq (10^{12}~\text{K})(r_g/r)\,;
\tag{4.7.5}
\end{equation}

and the synchrotron emissivity will retain its Newtonian value (4.4.4).  However,
the radiation which reaches infinity will be sharply cut off as the gas element
passes through the horizon.

\begin{equation}
\left(\text{fraction of emitted radiation that}\atop
\text{reaches } r = \infty \text{ from infalling gas at radius } r_f\right)
\simeq \frac{27}{64}\left(1 - \frac{r_g}{r_f}\right)\,.
\tag{4.7.6}
\end{equation}

In previous sections we have taken rough account of this cutoff and of the
accompanying redshift by retaining the Newtonian luminosity down to $r = 2r_g$.

In the case of a rotating (Kerr) hole, near and inside the ergosphere the
dragging of inertial frames will swing the infalling gas into orbital rotation
about the hole.  As the gas approaches the horizon, its angular velocity as seen
from infinity must approach the angular velocity of the horizon,

\begin{equation}
\Omega \to \Omega_{\text{horizon}} = \frac{a}{r_+^2 + a^2}
= \frac{c^3}{2GM}\,\frac{1}{(\text{sec})}\left(\frac{10^5}{M_*}\right)
\quad\text{(geometrized units, } c = G = 1\text{)}
\tag{4.7.7}
\end{equation}

for ``maximally rotating hole'', $a = M$.

Here $a$ is the hole's angular momentum per unit mass.  Any extra-luminous hot
spot in the gas (e.g.\ a region where magnetic field lines are reconnecting; cf.\
\S2.9) must orbit the hole with this angular velocity as it falls inward.  The
brightness of the hole as seen at Earth will be modulated with a period
$P \simeq 4\pi r_g/\Omega \gtrsim (10^{-4}~\text{sec})(M/M_\odot)$ as a result of the Doppler shift of the spot's
light and the bending of its rays.  (A factor of $2\pi/\Omega$ is the orbital period; the
other factor of $2\pi/\Omega$ is the light travel time between the radius at which the
spot is located at the beginning of one period and the radius to which it has
fallen at the end of the period.)  Thus, one might expect quasiperiodic
fluctuations in the light from a black hole living alone in the interstellar medium.

\bibliography{agent_15}

\end{document}
